\begin{proof}[Proof of \cref{obj:lem:amplification-dual-norm}]
Rewrite the unbiasedness equations in the single moment form
\[
\sum_{j=0}^k w_jp_j^\ell=
\begin{cases}
0,&\ell=0,\\
1,&1\le \ell\le\beta,
\end{cases}
\qquad 0\le\ell\le\beta .
\]
This is equivalent to \(w\in W_\beta(p)\): the case \(\ell=0\) is the zeroth-moment equation in \cref{obj:def:unbiased-weight-set}, while each positive \(\ell\le\beta\) is one of its remaining moment equations. % lean: key
Consequently the set of attainable \(\ell^1\) norms
\[
\left\{\sum_{j=0}^k |w_j|:w\in W_\beta(p)\right\}
\]
is exactly the primal norm set for this moment system, and the assumed nonemptiness of \(W_\beta(p)\) gives primal feasibility. % lean: hprimal_set

Finite-dimensional \(\ell^1/\ell^\infty\) duality for the endpoint functional \(r\mapsto r(1)-r(0)\), restricted to polynomials of degree at most \(\beta\), therefore gives
\[
\inf_{w\in W_\beta(p)}\sum_{j=0}^k |w_j|
=
\sup\left\{
|r(1)-r(0)|:
\deg(r)\le\beta,\ \max_{0\le j\le k}|r(p_j)|\le1
\right\}.
\] % lean: l1_repr_eq_sup_dual_of_momentSol_nonempty

For the squared identity, let
\[
D=
\sup\left\{
|r(1)-r(0)|:
\deg(r)\le\beta,\ \max_{0\le j\le k}|r(p_j)|\le1
\right\}.
\]
The zero polynomial is feasible on the dual side, so \(D\ge0\). % lean: dual_nonneg_of_momentSol_nonempty
Let
\[
\mathcal A=\left\{\left(\sum_{j=0}^k |w_j|\right)^2:w\in W_\beta(p)\right\}.
\]
This set is nonempty and bounded below by \(0\). For every feasible \(w\), the first duality equality gives \(D\le \sum_j |w_j|\), hence \(D^2\le(\sum_j |w_j|)^2\); taking the infimum over \(\mathcal A\) gives \(D^2\le\inf\mathcal A\). % lean: hAmp_lower
Conversely, strong duality supplies a feasible \(w\in W_\beta(p)\) with \(\sum_j |w_j|\le D\), and therefore
\[
\inf\mathcal A\le \left(\sum_j|w_j|\right)^2\le D^2.
\] % lean: exists_moment_le_dual_of_momentSol_nonempty
Thus \(\inf\mathcal A=D^2\). Since \cref{obj:def:amplification-criterion} defines \(A_\beta(p)\) as precisely this infimum, the claimed squared formula follows. % lean: amplification_dual_norm
\end{proof}
